\chapter{度量与 Levi--Civita 联络}

一般 仿射 联络 与长度、角度没有必然关系。若希望几何中存在“距离的无穷小版本”，必须额外指定度量；若进一步要求平行移动保持这个度量且没有挠率，联络就不再任意。Levi--Civita 定理说明：一个非退化对称度量唯一决定满足这两个条件的联络。本章先证明这一唯一性与存在性，再由同一联络定义测地线以及 Riemann、Ricci 与标量曲率，为后面的 Killing 对称、Hodge 理论和几何演化建立共同母结构。

\section{度量与正交标架}

Riemann 度量是在每点给定正定对称双线性型
\[
g_p:T_pM\times T_pM\to\RR
\]
且随 $p$ 光滑变化。若只要求非退化而允许固定符号差，则称 伪 Riemann 度量。局部坐标中
\begin{equation*}g=g_{ij}\,\dd x^i\otimes\dd x^j,
\qquad
g_{ij}=g_{ji},
\end{equation*}
逆矩阵记为 $g^{ij}$。

度量给出向量与一形式之间的自然同构
\begin{equation*}X^\flat:=g(X,\cdot),
\qquad
\alpha^\sharp\ \text{由}\ g(\alpha^\sharp,\cdot)=\alpha(\cdot)\ \text{定义}.
\end{equation*}
这是第一次允许“升降指标”。在此前各章中，上、下指标只表示张量槽位或标架分量，不能借助一个尚未定义的度量互换。

若标架 $\{e_a\}$ 满足
\[
g(e_a,e_b)=\eta_{ab}
\]
且 $\eta_{ab}$ 为固定对角矩阵，则称为局部正交归一标架。Riemann 情形 $\eta=I$；伪 Riemann 情形 $\eta$ 包含正负号。正交标架并不等于坐标标架：前者是度量条件，后者是 Lie 括号条件。

\section{Levi--Civita 定理与 Koszul 公式}

联络与度量相容指
\begin{equation*}\nabla g=0,
\end{equation*}
也就是
\[
X[g(Y,Z)]
=g(\nabla_XY,Z)+g(Y,\nabla_XZ).
\]
无挠条件则是 $T=0$，即
\[
\nabla_XY-\nabla_YX=[X,Y].
\]
两个条件联合起来足以完全解出 $\nabla_XY$。

\begin{theorem}[Levi--Civita 唯一性与存在性]\label{thm:math-dg-levi-civita}
给定任意 Riemann 或 伪 Riemann 度量 $g$，存在唯一 仿射 联络 $\nabla$ 同时满足
$\nabla g=0$ 与 $T=0$。它称为 $g$ 的 Levi--Civita 联络。
\end{theorem}

\begin{derivation}[从两条几何条件推出 Koszul 公式]
度量相容分别给出
\begin{align*}
Xg(Y,Z)&=g(\nabla_XY,Z)+g(Y,\nabla_XZ),\\
Yg(Z,X)&=g(\nabla_YZ,X)+g(Z,\nabla_YX),\\
Zg(X,Y)&=g(\nabla_ZX,Y)+g(X,\nabla_ZY).
\end{align*}
取前两式之和减第三式，得到
\begin{align*}
&Xg(Y,Z)+Yg(Z,X)-Zg(X,Y)\\
={}&g(\nabla_XY,Z)+g(Y,\nabla_XZ)
+g(\nabla_YZ,X)+g(Z,\nabla_YX)\\
&-g(\nabla_ZX,Y)-g(X,\nabla_ZY).
\end{align*}
现在逐次使用无挠关系
\[
\nabla_XZ-\nabla_ZX=[X,Z],
\quad
\nabla_YX-\nabla_XY=[Y,X],
\quad
\nabla_YZ-\nabla_ZY=[Y,Z].
\]
把所有出现的联络项整理到 $g(\nabla_XY,Z)$，得到
\begin{align*}
2g(\nabla_XY,Z)
={}&Xg(Y,Z)+Yg(Z,X)-Zg(X,Y)\\
&+g([X,Y],Z)-g([Y,Z],X)+g([Z,X],Y).
\end{align*}
因此
\begin{equation}\label{eq:math-dg-koszul-formula}
\boxed{
\begin{aligned}
2g(\nabla_XY,Z)
={}&Xg(Y,Z)+Yg(Z,X)-Zg(X,Y)\\
&+g([X,Y],Z)-g([Y,Z],X)+g([Z,X],Y).
\end{aligned}}
\end{equation}
右端只含已知的 $g$、向量场与 Lie 括号。由于 $g$ 非退化，\eqnrefs{eq:math-dg-koszul-formula}对每个 $X,Y$ 唯一确定 $\nabla_XY$，故唯一性成立。

反过来，以\eqnrefs{eq:math-dg-koszul-formula}定义 $\nabla_XY$。右端对 $Z$ 是函数线性的，故由非退化性确实唯一对应一个向量场。直接检查 $X,Y$ 上的线性与 Leibniz 规则可知它是仿射联络。将 $X,Y$ 互换并相减，右端化为 $2g([X,Y],Z)$，故 $T=0$；把 Koszul 公式适当置换后相加，则恢复
\[
Xg(Y,Z)=g(\nabla_XY,Z)+g(Y,\nabla_XZ),
\]
故 $\nabla g=0$。存在性也得到证明。

Koszul 公式是唯一性定理：给定度量 $g$，Levi--Civita 联络（无挠+度量相容）被 $g$ 唯一决定，无需选择坐标或标架。实例：欧氏度量 $g=\delta_{ij}\dd x^i\dd x^j$ 下，所有 $\partial_k g_{ij}=0$，Koszul 公式给出 $\Gamma^k_{ij}=0$（平直联络）。球面度量 $g=R^2(\dd\theta^2+\sin^2\theta\,\dd\varphi^2)$ 下，代入得 $\Gamma^\theta_{\varphi\varphi}=-\sin\theta\cos\theta$、$\Gamma^\varphi_{\theta\varphi}=\cot\theta$，其余为零——这就是球面测地线方程的联络系数。
\end{derivation}

在坐标基 $\partial_i$ 中 Lie 括号消失，把
$\nabla_{\partial_i}\partial_j=\Gamma^k{}_{ij}\partial_k$ 代入 Koszul 公式可得
\begin{equation*}\boxed{
\Gamma^k{}_{ij}
=\frac12g^{k\ell}
\left(
\partial_i g_{j\ell}
+\partial_j g_{i\ell}
-\partial_\ell g_{ij}
\right).}
\end{equation*}
这只是 Levi--Civita 联络在坐标基中的表达，不应反过来把 Christoffel 符号当作联络的坐标无关定义。

在一般标架中，Koszul 公式自动包含结构函数。令
\[
g_{ab}=g(e_a,e_b),
\qquad
[e_a,e_b]=C^c{}_{ab}e_c,
\qquad
\Gamma_{abc}:=g(\nabla_{e_a}e_b,e_c).
\]
则
\begin{equation*}\begin{aligned}
2\Gamma_{abc}
={}&e_a(g_{bc})+e_b(g_{ca})-e_c(g_{ab})\\
&+C^d{}_{ab}g_{dc}-C^d{}_{bc}g_{da}+C^d{}_{ca}g_{db}.
\end{aligned}
\end{equation*}
若标架正交归一，$g_{ab}=\eta_{ab}$ 为常数，前三个导数项消失，联络完全由非完整系数决定。

对正交归一标架，度量相容还给出
\begin{equation}\label{eq:math-dg-orthonormal-connection-skew}
\omega_{ab}+\omega_{ba}=0,
\qquad
\omega_{ab}:=\eta_{ac}\omega^c{}_b.
\end{equation}
于是第一结构方程在 Levi--Civita 情形化为
\[
\dd\theta^a+\omega^a{}_b\wedge\theta^b=0,
\]
常常可以直接从 $\dd\theta^a$ 解出联络一形式。

\section{测地线与曲率收缩}

给定 Levi--Civita 联络后，曲线 $\gamma(s)$ 的协变加速度为 $\nabla_{\dot\gamma}\dot\gamma$。若
\begin{equation*}\nabla_{\dot\gamma}\dot\gamma=0,
\end{equation*}
则称 $\gamma$ 为仿射参数测地线。坐标中它成为
\[
\frac{\dd^2x^k}{\dd s^2}
+\Gamma^k{}_{ij}
\frac{\dd x^i}{\dd s}
\frac{\dd x^j}{\dd s}=0.
\]
这个方程来自“切向量沿自身平行”，并不依赖把流形嵌入更大的 Euclidean 空间。

采用上一章的曲率约定
\[
R(X,Y)Z
=\nabla_X\nabla_YZ-\nabla_Y\nabla_XZ-\nabla_{[X,Y]}Z.
\]
对 Levi--Civita 联络，四阶协变张量
\[
R(X,Y,Z,W):=g(R(X,Y)Z,W)
\]
满足
\begin{align*}
R(X,Y,Z,W)&=-R(Y,X,Z,W),\\
R(X,Y,Z,W)&=-R(X,Y,W,Z),\\
R(X,Y,Z,W)&=R(Z,W,X,Y),
\end{align*}
并在无挠条件下满足代数第一 Bianchi 恒等式
\[
R(X,Y)Z+R(Y,Z)X+R(Z,X)Y=0.
\]
这些对称性大幅减少曲率的独立分量。

Ricci 张量定义为曲率算子的迹：若 $\{e_a\}$ 是局部正交标架，符号记为 $\varepsilon_a=g(e_a,e_a)=\pm1$，则
\begin{equation}\label{eq:math-dg-ricci-definition}
\operatorname{Ric}(Y,Z)
:=\sum_a\varepsilon_a\,g(R(e_a,Y)Z,e_a).
\end{equation}
标量曲率再作一次迹：
\begin{equation*}\mathcal R:=\operatorname{tr}_g\operatorname{Ric}.
\end{equation*}
Riemann 曲率保留二维方向之间的完整信息，Ricci 曲率是其一次收缩，标量曲率则只保留总迹。后面 Ricci 流 的母方程会直接使用\eqnrefs{eq:math-dg-ricci-definition}。
Ricci 与标量曲率都是收缩后的信息；若要保留一张二维切平面本身的弯曲强度，应先定义截面曲率。对点 $p\in M$ 处由线性无关向量 $u,v$ 张成的二维平面 $\sigma\subset T_pM$，定义
\begin{equation*}K(\sigma)
:=\frac{R(u,v,v,u)}{g(u,u)g(v,v)-g(u,v)^2}.
\end{equation*}
分母正是 Gram 行列式，因此在 Riemann 情形严格为正，并且右端只依赖二维子空间 $\sigma$，与所选基 $u,v$ 无关。若 $u,v$ 正交归一，便简化为
\[
K(u,v)=R(u,v,v,u).
\]
于是曲率张量可以理解为“所有二维方向上的无穷小偏离平坦程度”的统一编码。

固定内积后，可把 Riemann 曲率等价看作 $\Lambda^2T_pM$ 上的自伴算子。定义 $\mathcal R_p:\Lambda^2T_pM\to\Lambda^2T_pM$ 使
\begin{equation*}\left\langle \mathcal R_p(u\wedge v),w\wedge z\right\rangle
:=R(u,v,z,w).
\end{equation*}
这里 $\Lambda^2T_pM$ 上的内积由 $g$ 诱导。曲率张量的配对对称性保证 $\mathcal R_p$ 自伴。于是
\[
K(\sigma)
=\frac{\langle \mathcal R_p\xi,\xi\rangle}{\langle\xi,\xi\rangle},
\qquad
\xi=u\wedge v,
\]
恰是曲率算子在 decomposable 双矢 上的 Rayleigh quotient。这个观点把 比较 几何 与前面谱理论的语言直接接起来：截面曲率上下界就是曲率算子在所有简单二向量方向上的二次型界。

Ricci 曲率可由截面曲率求和恢复。若 $v$ 是单位向量，并取正交归一标架
\[
e_1=v,e_2,\ldots,e_n,
\]
则
\begin{equation*}\operatorname{Ric}(v,v)
=\sum_{a=2}^{n}K(v,e_a).
\end{equation*}
所以 Ricci 曲率不是某个单独二维平面的曲率，而是所有包含 $v$ 的正交二维平面的平均型总和；标量曲率则进一步满足
\begin{equation*}\mathcal R
=2\sum_{1\le a<b\le n}K(e_a,e_b).
\end{equation*}
这解释了为什么 Bonnet--Myers 只需 Ricci 下界，而 Rauch 比较通常需要逐截面曲率界：两者控制的是不同强度的聚焦信息。

若所有点、所有二维截面都有同一常数曲率 $K$，则曲率张量被唯一逼成
\begin{equation}\label{eq:math-dg-space-form-curvature}
R(X,Y)Z
=K\bigl(g(Y,Z)X-g(X,Z)Y\bigr).
\end{equation}
反之，\eqnrefs{eq:math-dg-space-form-curvature} 给出常截面曲率 $K$。Euclidean 空间、半径 $a$ 的圆球以及常曲率双曲空间分别对应 $K=0$、$K=a^{-2}$ 与 $K=-a^{-2}$。它们不是三个互不相干的例子，而是所有 比较定理 的标准参照模型。

\begin{derivation}[由截面曲率恢复代数曲率张量]
设 $R$ 与 $\widetilde R$ 是两个具有 Riemann 曲率全部代数对称性的四阶张量，并且它们给出相同截面曲率。令
\begin{equation*}
A:=R-\widetilde R.
\end{equation*}
则 $A$ 仍具有 Riemann 代数对称性，且对任意单位正交 $u,v$，
\begin{equation*}
A(u,v,v,u)=R(u,v,v,u)-\widetilde R(u,v,v,u)=0.
\end{equation*}
由 $A$ 在每个槽位关于度量的双线性，对任意 $u,v$（不必单位正交）也有
\begin{equation*}
A(u,v,v,u)=0.
\end{equation*}
我们证明这已迫使 $A=0$。先把 $v$ 替换为 $v+w$，利用双线性展开和 $A$ 在最后两个槽位的反对称性，
\begin{align*}
0&=A(u,v+w,v+w,u)\\
&=A(u,v,v,u)+A(u,v,w,u)+A(u,w,v,u)+A(u,w,w,u)\\
&=A(u,v,w,u)+A(u,w,v,u),
\end{align*}
即
\begin{equation*}
A(u,v,w,u)+A(u,w,v,u)=0.
\end{equation*}
再把 $u$ 替换为 $u+z$，展开并利用上式及配对对称性 $A(u,v,w,z)=A(w,z,u,v)$，得到
\begin{align*}
0&=A(u+z,v,w,u+z)\\
&=A(u,v,w,u)+A(u,v,w,z)+A(z,v,w,u)+A(z,v,w,z)\\
&=A(u,v,w,z)+A(z,v,w,u).
\end{align*}
由第一 Bianchi 恒等式 $A(u,v,w,z)+A(v,w,u,z)+A(w,u,v,z)=0$ 与上式反复换位，可把任意 $A(u,v,w,z)$ 极化成若干个 $A(a,b,b,a)$ 的线性组合。一个显式极化形式是先定义
\begin{equation*}
Q(a,b):=A(a,b,b,a),
\end{equation*}
则 $Q$ 是 $(a,b)$ 上的双线性型；连续对 $a$ 与 $b$ 做二次型极化，
\begin{align*}
A(u,v,w,z)
&=\frac14\Bigl\{Q(u+w,v+z)-Q(u-w,v+z)\\
&\quad-Q(u+w,v-z)+Q(u-w,v-z)\Bigr\}
\end{align*}
（具体系数由上述代数对称性确定），所有项都由 $Q$ 决定。由于假设 $Q\equiv0$，于是 $A\equiv0$。

因此在 Riemann 情形，知道每一点所有二维截面的 $K(\sigma)$，就等价于知道完整的 Riemann 曲率张量。所谓"截面曲率是曲率的基本标量"并不是丢掉信息后的口号，而是依靠曲率张量特殊代数对称性的可逆编码。
\end{derivation}

\section{指数映射、Jacobi 场与完备性}

测地线方程是二阶常微分方程，因此对每个 $p\in M$ 和初速度 $v\in T_pM$，局部存在唯一测地线 $\gamma_v$ 满足
\[
\gamma_v(0)=p,
\qquad
\dot\gamma_v(0)=v.
\]
只要该测地线至少存在到参数 $1$，定义指数映射
\begin{equation*}\exp_p(v):=\gamma_v(1).
\end{equation*}
由测地线方程对初值的光滑依赖，$\exp_p$ 在零向量附近光滑，并且
\[
(\dd\exp_p)_0=\id(T_pM).
\]
所以逆函数定理说明 $\exp_p$ 在零向量附近给出一组 正规 坐标。这里“正规”不是额外坐标假设，而是由所有从 $p$ 发出的测地线同时构造出的坐标。

相邻测地线之间的一阶分离由 Jacobi 场控制。设 $F(s,r)$ 是测地线变分，即对每个固定 $r$，$s\mapsto F(s,r)$ 都是测地线。沿中心测地线 $\gamma(s)=F(s,0)$ 定义
\[
T:=\partial_sF\big|_{r=0},
\qquad
J:=\partial_rF\big|_{r=0}.
\]
则 $[T,J]=0$，且 $J$ 满足
\begin{equation}\label{eq:math-dg-jacobi-equation}
\boxed{\nabla_T\nabla_TJ+R(J,T)T=0.}
\end{equation}

\begin{derivation}[由测地线变分推出 Jacobi 方程，即\eqnrefs{eq:math-dg-jacobi-equation}]
$F(s,r)$ 是测地线变分：对每个固定 $r$，$s\mapsto F(s,r)$ 都是测地线。沿中心测地线 $\gamma(s)=F(s,0)$ 定义切场 $T=\partial_sF|_{r=0}$ 和变分场 $J=\partial_rF|_{r=0}$。因混合偏导交换，$[T,J]=0$。

由于每条 $r=\mathrm{const.}$ 曲线都是测地线，
\begin{equation*}
 \nabla_TT=0.
\end{equation*}
这是测地线的内蕴定义，等价于 $\dd^2 x^k/\dd s^2+\Gamma^k{}_{ij}\dot x^i\dot x^j=0$。

Levi--Civita 联络无挠，故 $\nabla_TJ-\nabla_JT=[T,J]=0$，即
\begin{equation*}
 \nabla_TJ=\nabla_JT.
\end{equation*}
于是
\begin{equation*}
 \nabla_T\nabla_TJ=\nabla_T\nabla_JT.
\end{equation*}
这一步把"分离场沿测地线的二阶协变导数"转化为"切场沿变分方向的协变导数"，是后续代入曲率定义的关键准备。

按本书曲率约定，
\begin{equation*}
 R(T,J)T
 =\nabla_T\nabla_JT-\nabla_J\nabla_TT-\nabla_{[T,J]}T.
\end{equation*}
后两项分别因 $\nabla_TT=0$（测地线方程）与 $[T,J]=0$（混合偏导交换）消失，故
\begin{equation*}
 R(T,J)T=\nabla_T\nabla_JT=\nabla_T\nabla_TJ.
\end{equation*}

由曲率算子的反对称性 $R(T,J)T=-R(J,T)T$，
\begin{equation*}
 \nabla_T\nabla_TJ=-R(J,T)T,
\end{equation*}
移项即得\eqnrefs{eq:math-dg-jacobi-equation}。

球面与双曲空间中的 Jacobi 场：
\begin{itemize}
\item 单位球面 $S^n$（截面曲率 $+1$）：沿大圆测地线取单位正交初始 Jacobi 场 $J(0)=0$、$\nabla_TJ(0)=e$，方程 $\nabla_T^2 J+J=0$，解
\begin{equation*}
 J(s)=\sin(s)\,e(s),
\end{equation*}
$s=\pi$ 处 $J(\pi)=0$，即对径点为共轭点。所有大圆在对径点处汇聚。
\item 双曲空间 $\HH^n$（截面曲率 $-1$）：方程 $\nabla_T^2 J-J=0$，解
\begin{equation*}
 J(s)=\sinh(s)\,e(s),
\end{equation*}
$\sinh(s)$ 在 $s>0$ 不为零，无共轭点。测地线指数发散，这是 Cartan--Hadamard 定理的典型实例。
\item $\RR^n$（截面曲率 $0$）：方程 $\nabla_T^2 J=0$，解 $J(s)=s\,e$，线性分离，无共轭点。
\end{itemize}
三种曲率符号对应汇聚（$\sin$）、发散（$\sinh$）、线性（$s$），是 Rauch 比较定理的微缩版。

Jacobi 方程描述相邻测地线的相对加速度：$\nabla_T\nabla_TJ$ 是分离加速度，$R(J,T)T$ 是曲率对分离的回复力。正曲率使测地线汇聚（共轭点，如引力聚焦），负曲率使测地线发散（如宇宙膨胀的混沌混合）。在广义相对论中，Jacobi 方程是测地偏离方程，描述相邻自由下落粒子的相对加速度——$R(J,T)T$ 正是潮汐力张量作用于分离矢量。这是引力场探测（如 LIGO 引力波探测）的理论基础：引力波通过改变曲率 $R$，使相邻测地线的分离 $J$ 产生可测振荡。
\end{derivation}

若存在非零 Jacobi 场满足 $J(0)=J(s_*)=0$，则 $\gamma(0)$ 与 $\gamma(s_*)$ 沿该测地线互为共轭点；等价地，$(\dd\exp_p)_{s_*v}$ 在相应方向失去可逆性。因而指数映射的局部退化、测地线族聚焦与曲率通过同一 Jacobi 方程联系起来。

完备性把这种局部理论延伸到全局。对连通 Riemann 流形，Hopf--Rinow 定理给出一组等价条件：度量空间 $(M,d_g)$ 完备；所有闭有界子集紧；任意测地线可延拓到全部实参数；任意两点之间存在实现距离的最短测地线。特别地，测地完备并不是“坐标能覆盖全局”，而是指数映射 $\exp_p$ 对每个 $p$ 都定义在整个 $T_pM$ 上。
指数映射不仅给出坐标，还把径向方向与角向方向正交分离。固定 $p\in M$，在 $T_pM$ 的零向量附近考虑
\[
\exp_p:T_pM\to M.
\]
对 $v\in T_pM$ 与任意 $w\in T_pM$，Gauss 引理断言
\begin{equation}\label{eq:math-dg-gauss-lemma}
 g_{\exp_p(v)}\!\left(
 (d\exp_p)_v v,
 (d\exp_p)_v w
 \right)
 =g_p(v,w).
\end{equation}
特别地，若 $w\perp v$，则径向测地线方向与由球面 $|v|=\mathrm{const.}$ 推出的角向方向正交。这一正交性是“测地极坐标”成立的真正原因。

\begin{derivation}[Gauss 引理由测地线能量守恒推出]
取从 $p$ 出发的二维测地线变分，沿径向参数 $t$ 求内积 $g(T,J)$ 的导数。利用测地线方程 $\nabla_TT=0$、无挠性 $\nabla_TJ=\nabla_JT$ 与速度平方沿测地线守恒，把导数化为 $\tfrac12 J[g(T,T)]$，再用速度守恒化简并沿径向 $t$ 积分得 Gauss 引理。

取二维变分
\begin{equation*}
F(t,s)=\exp_p\bigl(t(v+s w)\bigr),
\qquad 0\le t\le1.
\end{equation*}
对固定 $s$，$t\mapsto F(t,s)$ 是测地线（指数映射的定义）。记
\begin{equation*}
T=\partial_tF,
\qquad
J=\partial_sF.
\end{equation*}
$T$ 是测地线切向，$J$ 是变分场。在 $t=0$ 处 $F(0,s)=p$ 不依赖 $s$，故 $J(0)=0$；在 $t=1$ 处 $F(1,s)=\exp_p(v+sw)$，$J(1)=(d\exp_p)_v(w)$。同时 $T(1)=(d\exp_p)_v(v)$。

由于 $[T,J]=0$（混合偏导交换）且 Levi--Civita 联络无挠，$\nabla_TJ=\nabla_JT$。

沿 $t$ 方向求
\begin{equation*}
\frac{\dd}{\dd t}g(T,J)
=g(\nabla_TT,J)+g(T,\nabla_TJ).
\end{equation*}
第一项因测地线方程 $\nabla_TT=0$ 消失。第二项用 $\nabla_TJ=\nabla_JT$：
\begin{equation*}
\frac{\dd}{\dd t}g(T,J)
=g(T,\nabla_JT)
=\frac12 J\bigl[g(T,T)\bigr].
\end{equation*}
最后等式用了度量相容性 $J[g(T,T)]=2g(\nabla_JT,T)$。

每条径向测地线速度平方恒定，并且初速度为 $v+s w$，故
\begin{equation*}
J[g(T,T)]
=\frac{\partial}{\partial s}g_p(v+s w,v+s w)
=2g_p(v+s w,w).
\end{equation*}
关键处用了"测地线速度的范数沿测地线守恒"——这是 $\nabla_TT=0$ 与度量相容性的直接推论：$\frac{\dd}{\dd t}g(T,T)=2g(\nabla_TT,T)=0$。

在 $s=0$ 处积分 $t\in[0,1]$：
\begin{equation*}
g(T(1),J(1))-g(T(0),J(0))
=\int_0^1 g_p(v,w)\,\dd t
=g_p(v,w).
\end{equation*}
又因 $J(0)=0$，左端第一项 $g(T(0),J(0))=0$，得到
\begin{equation*}
g(T(1),J(1))=g_p(v,w),
\end{equation*}
即 \eqnrefs{eq:math-dg-gauss-lemma}。

几何意义：
\begin{itemize}
\item $w\perp v$ 时 $g_p(v,w)=0$，故径向方向 $(d\exp_p)_v v$ 与角向方向 $(d\exp_p)_v w$ 在 $\exp_p(v)$ 处正交。这是"测地极坐标 $g=\dd r^2+r^2 h(r,\theta)$"成立的真正原因。
\item $\RR^n$：$\exp_p(v)=p+v$，$(d\exp_p)_v=\mathrm{id}$，Gauss 引理退化为 $g(v,w)=g(v,w)$，平凡成立。
\item $S^n$：$\exp_p(v)$ 把 $T_pS^n$ 中向量映到球面，Gauss 引理保证球面测地极坐标 $g=\dd r^2+\sin^2(r)\,d\Omega^2$ 中径向与角向正交。
\end{itemize}

Gauss 引理是 Rauch 比较定理的局部基石：它保证指数映射在径向方向不改变长度，因此比较定理只需控制角向方向的伸缩。在广义相对论中，Gauss 引理保证正规坐标中径向光锥的几何性质，是因果结构分析的基础工具。
\end{derivation}

在 $p$ 的 正规 坐标 $(x^1,\ldots,x^n)$ 中，坐标线 $t\mapsto tx$ 是从 $p$ 发出的测地线。由测地线方程在 $t=0$ 对任意初速度成立，得到
\begin{equation*}g_{ij}(p)=\delta_{ij},
\qquad
\Gamma^k{}_{ij}(p)=0,
\qquad
\partial_k g_{ij}(p)=0.
\end{equation*}
因此 正规 坐标 消去的是度量的一阶变化，而不是曲率。真正的几何信息从二阶 Taylor 系数开始出现。与本书曲率约定相容的展开为
\begin{equation*}g_{ij}(x)
 =\delta_{ij}
 -\frac13 R_{ikj\ell}(p)x^k x^\ell
 +O(|x|^3),
\end{equation*}
从而
\begin{equation}\label{eq:math-dg-normal-volume-expansion}
 \sqrt{\det g(x)}
 =1-\frac16\operatorname{Ric}_{k\ell}(p)x^k x^\ell
 +O(|x|^3).
\end{equation}
这两式给出了“曲率是度量偏离 Euclidean 形式的二阶障碍”的严格含义。Riemann 张量控制各方向的二阶度量畸变，Ricci 张量则控制小体积元的二阶畸变。

由 \eqnrefs{eq:math-dg-normal-volume-expansion} 对半径 $r$ 的小测地球积分，可得
\begin{equation*}\operatorname{Vol}_g(B_r(p))
=\omega_n r^n
\left[
1-\frac{\mathcal R(p)}{6(n+2)}r^2+O(r^3)
\right],
\end{equation*}
其中 $\omega_n$ 是 Euclidean 单位 $n$-球体积。正标量曲率在最低非平凡阶压缩小球体积，负标量曲率则扩张它。
正规 坐标 只能在指数映射保持一一且微分可逆的区域内工作。固定 $p$，对单位向量 $v\in T_pM$ 定义 割 time
\begin{equation*}c(v):=\sup\left\{t>0:
 d\bigl(p,\exp_p(sv)\bigr)=s
 \text{ for every }0\le s\le t
 \right\}.
\end{equation*}
若 $c(v)<\infty$，则 $\exp_p(c(v)v)$ 称为沿 $v$ 的 割 point；所有这类点组成 $\operatorname{Cut}(p)$。单射半径 定义为
\begin{equation*}\operatorname{inj}(p):=\sup\{r>0:
 \exp_p|_{B_r(0)}\text{ is a diffeomorphism onto its image}\}.
\end{equation*}
在完备 Riemann 流形上，$\operatorname{inj}(p)$ 等于从 $p$ 出发最早发生“共轭退化”或“不同最短测地线竞争”时的半径。前者由 $d\exp_p$ 奇异造成，后者即使 $d\exp_p$ 仍可逆也会破坏全局一一性。

在
\[
M\setminus\bigl(\{p\}\cup\operatorname{Cut}(p)\bigr)
\]
上，距离函数
\[
r_p(q):=d(p,q)
\]
光滑，并且
\begin{equation*}|\nabla r_p|=1,
\qquad
\nabla r_p(q)=\dot\gamma_{pq}(r_p(q)),
\end{equation*}
其中 $\gamma_{pq}$ 是唯一的单位速最短测地线。第一式是 Riemann eikonal 方程。它说明距离函数虽是全局 Lipschitz 对象，但在 割迹 之外服从一阶 Hamilton--Jacobi 方程；奇性恰在最短测地线开始多值或指数映射退化的地方产生。

\begin{theorem}[最短性与 割/conjugate 障碍]\label{thm:math-dg-minimizing-cut-conjugate}
设 $\gamma:[0,L]\to M$ 为单位速测地线。若在 $(0,L]$ 内出现与 $\gamma(0)$ 共轭的点，则 $\gamma$ 在越过第一个共轭点后不再是局部最短曲线。即使没有共轭点，若 $\gamma(L)\in\operatorname{Cut}(\gamma(0))$，也可能由于存在另一条同长最短测地线而失去唯一最短性。
\end{theorem}

这一区分非常重要：Jacobi 场和 指标 形式 检测的是指数映射的微分退化，是“局部二阶”障碍；割迹 还记录不同测地线之间的全局竞争。圆球上南极同时是北极的共轭点与 割 point；平坦圆柱上某些 割 point 却可由拓扑导致，而局部曲率仍为零。

测地线还具有一个等价的变分刻画。设 $\gamma:[0,L]\to M$ 为分段光滑曲线，能量泛函定义为
\begin{equation*}E[\gamma]
:=\frac12\int_0^L g(\dot\gamma,\dot\gamma)\,\dd s.\end{equation*}
对固定端点变分 $F(s,r)$，记变分场
\[
V(s):=\partial_rF(s,r)|_{r=0},
\qquad V(0)=V(L)=0.
\]
一阶变分公式为
\begin{equation}
\delta E[V]
=-\int_0^L g(V,\nabla_TT)\,\dd s,
\qquad T=\dot\gamma.
\label{eq:math-dg-energy-first-variation}
\end{equation}
所以固定端点下的能量临界曲线恰好是仿射参数测地线。这个结论解释了为什么测地线方程同时具有“平行运输”与“变分极值”两种来源。

对一条测地线再做一次变分，得到 指标 形式
\begin{equation}
I(V,W)
:=\int_0^L
\left[
 g(\nabla_TV,\nabla_TW)
 -g(R(V,T)T,W)
\right]\dd s,
\label{eq:math-dg-index-form}
\end{equation}
且固定端点时
\[
\delta^2E[V,W]=I(V,W).
\]
因此曲率不是只改变“邻近测地线是否聚焦”；它直接进入能量的 Hessian，决定测地线在变分意义下是否稳定。

\begin{derivation}[从能量变分到指标形式]
对二维变分 $F(s,r)$ 记
\[
T=\partial_sF,
\qquad
V=\partial_rF.
\]
坐标向量场满足 $[T,V]=0$，Levi--Civita 联络无挠，因此
\[
\nabla_VT=\nabla_TV.
\]
对能量求一次导数：
\[
\begin{aligned}
\frac{\dd E}{\dd r}
&=\int_0^L g(\nabla_VT,T)\,\dd s\\
&=\int_0^L g(\nabla_TV,T)\,\dd s\\
&=\bigl[g(V,T)\bigr]_0^L
-\int_0^L g(V,\nabla_TT)\,\dd s.
\end{aligned}
\]
固定端点使边界项消失，得到\eqnrefs{eq:math-dg-energy-first-variation}。

现在假设中心曲线 $r=0$ 是测地线。对一阶变分式关于 $r$ 再求导，并在 $r=0$ 取值。由于 $\nabla_TT=0$，只有 $\nabla_V\nabla_TT$ 的变化留下：
\[
\delta^2E[V,V]
=-\int_0^L g\bigl(V,\nabla_V\nabla_TT\bigr)\,\dd s.
\]
曲率恒等式与 $[V,T]=0$ 给出
\[
\nabla_V\nabla_TT
=\nabla_T\nabla_VT+R(V,T)T
=\nabla_T\nabla_TV+R(V,T)T.
\]
于是
\[
\delta^2E[V,V]
=-\int_0^L g(V,\nabla_T\nabla_TV)\,\dd s
-\int_0^L g(R(V,T)T,V)\,\dd s.
\]
对第一项分部积分，固定端点给出
\[
-\int_0^L g(V,\nabla_T\nabla_TV)\,\dd s
=\int_0^L|\nabla_TV|^2\,\dd s.
\]
故
\begin{equation}
\boxed{
\delta^2E[V,V]
=\int_0^L
\left(
|\nabla_TV|^2-g(R(V,T)T,V)
\right)\dd s.}
\label{eq:math-dg-energy-second-variation}
\end{equation}
极化后即得到\eqnrefs{eq:math-dg-index-form}。

指标形式 $I(V,V)=\int(|\nabla_TV|^2-g(R(V,T)T,V))\,\dd s$ 的两项竞争：协变导数项 $|\nabla_TV|^2\ge0$ 使测地线稳定，曲率项 $-g(R(V,T)T,V)$ 在正曲率方向为负使测地线不稳定。若存在 $V\neq0$ 使 $I(V,V)<0$，则测地线不是能量极小——存在更短的邻近曲线。使 $I(J,J)=0$ 的 Jacobi 场 $J$ 标志共轭点：测地线过了共轭点后不再极小。

在 $\mathbb S^n(R)$ 上取赤道测地线 $\gamma(s)$（$0\le s\le\pi R$）。沿 $\gamma$ 取法向变分场 $V=f(s)\vb e_\perp$，其中 $\vb e_\perp$ 平行单位法向量。球面截面曲率 $K=1/R^2$，故
\[
 I(V,V)=\int_0^L\!\left(|f'|^2-\frac{1}{R^2}|f|^2\right)\dd s.
\]
取 $f(s)=\sin(s/R)$（满足 $f(0)=f(\pi R)=0$），则 $I=0$——赤道在 $s=\pi R$（对跖点）处有共轭点。对 $L>\pi R$，取 $f$ 使 $I<0$，测地线不再极小。这就是球面上两点间测地线只在距离 $<\pi R$ 时最短的变分解释。
\end{derivation}

能量的第二变分公式由此给出：固定端点时 $\delta^2 E=\int_0^L(|\nabla_T V|^2-g(V,\nabla_T\nabla_TV))\dd s$（\eqnrefs{eq:math-dg-energy-second-variation}）。

Jacobi 场正是 指标 形式 的零模。若 $J(0)=J(L)=0$ 并满足 Jacobi 方程，则对任意固定端点变分场 $W$，分部积分得到
\[
I(J,W)
=-\int_0^L
 g\bigl(\nabla_T\nabla_TJ+R(J,T)T,W\bigr)\,\dd s=0.
\]
因此共轭点出现时，能量 Hessian 退化。更进一步，Morse 指标定理 说明一条测地线段的 指标 形式 负方向数等于区间内共轭点的重数总和。于是“共轭点前局部极小、越过共轭点后出现下降方向”不是几何图像上的猜测，而是二阶变分的谱性质。这里说的是固定端点的局部极小性；全局最短性还会受到 割迹 和不同测地线竞争的影响。


曲率的真正威力来自比较：Jacobi 方程把一个张量不等式转化成测地线会聚速度的不等式。先看常截面曲率 $K$ 的模型空间。若 $J\perp T$ 且沿测地线选平行单位场 $E$ 写成
\[
J(r)=j(r)E(r),
\]
则 Jacobi 方程化为标量 ODE
\[
j''+Kj=0.
\]
满足 $j(0)=0$、$j'(0)=1$ 的模型解统一写成
\begin{equation}
S_K(r)
:=
\begin{cases}
\dfrac{\sin(\sqrt K\,r)}{\sqrt K},&K>0,\\[7pt]
r,&K=0,\\[5pt]
\dfrac{\sinh(\sqrt{-K}\,r)}{\sqrt{-K}},&K<0.
\end{cases}
\label{eq:math-dg-model-jacobi}
\end{equation}
因此正曲率使径向测地线比 Euclidean 情形更快聚焦，负曲率则使其更快分离。$S_K$ 并不只是三个特例的记忆公式；它是 比较 几何 中所有 radial estimate 的母函数。

Rauch 比较定理把这一模型推广到变曲率情形。粗略地说，若两条单位速测地线上的对应二维截面满足
\[
K_1\le K_2,
\]
并以相同初始横向速度发出 Jacobi 场，则在第一个共轭点以前，较小曲率一侧的 $|J|$ 不小于较大曲率一侧。也就是说，``更正的曲率 $\Rightarrow$ 更强的聚焦''可以被严格比较，而不需要显式求解 Jacobi 方程。

这种比较立刻给出两个方向相反的全局结论。若完整、单连通 Riemann 流形满足
\[
K\le0,
\]
则 Cartan--Hadamard 定理断言任意 $p\in M$ 的指数映射
\[
\exp_p:T_pM\to M
\]
是覆盖映射；在单连通情形进一步是全局微分同胚。于是非正曲率消除了共轭点和测地线的局部聚焦障碍，$T_pM$ 本身给出全局坐标模型。

正 Ricci 曲率则产生相反的紧致化效应。设 $\dim M=n$ 且
\begin{equation*}\operatorname{Ric}(v,v)
\ge (n-1)k\,g(v,v),
\qquad k>0.\end{equation*}
Bonnet--Myers 定理给出
\begin{equation}
\boxed{
\operatorname{diam}(M)
\le\frac{\pi}{\sqrt k}
},
\label{eq:math-dg-bonnet-myers-diameter}
\end{equation}
并推出完备流形 $M$ 紧致、基本群有限。关键点是：这里不需要每个截面曲率都有正下界，只需所有与测地线方向相交的截面曲率之和，也就是 Ricci 曲率，有统一正下界。

\begin{derivation}[Bonnet--Myers 如何从指标形式推出直径上界]
反设存在两点距离
\[
L>\frac{\pi}{\sqrt k}.
\]
由 Hopf--Rinow，存在连接它们的单位速最短测地线
\[
\gamma:[0,L]\to M.
\]
沿 $\gamma$ 取平行正交标架
\[
E_1,\ldots,E_{n-1},E_n=T:=\dot\gamma.
\]
对 $i=1,\ldots,n-1$ 定义固定端点变分场
\[
V_i(t)
=\sin\!\left(\frac{\pi t}{L}\right)E_i(t).
\]
由于 $\gamma$ 是最短测地线，在未越过割点的区间上，指标形式必须非负：
\[
I(V_i,V_i)\ge0.
\]
直接计算
\[
|\nabla_TV_i|^2
=\frac{\pi^2}{L^2}
\cos^2\!\left(\frac{\pi t}{L}\right),
\]
以及
\[
g(R(V_i,T)T,V_i)
=\sin^2\!\left(\frac{\pi t}{L}\right)
K(T,E_i).
\]
对 $i$ 求和，并使用
\[
\sum_{i=1}^{n-1}K(T,E_i)
=\operatorname{Ric}(T,T)
\ge(n-1)k,
\]
得到
\[
\begin{aligned}
0
&\le\sum_{i=1}^{n-1}I(V_i,V_i)\\
&\le(n-1)
\int_0^L
\left[
\frac{\pi^2}{L^2}
\cos^2\!\left(\frac{\pi t}{L}\right)
-k\sin^2\!\left(\frac{\pi t}{L}\right)
\right]\dd t.
\end{aligned}
\]
利用
\[
\int_0^L\sin^2\!\left(\frac{\pi t}{L}\right)\dd t
=
\int_0^L\cos^2\!\left(\frac{\pi t}{L}\right)\dd t
=\frac L2,
\]
可得
\[
0
\le
\frac{(n-1)L}{2}
\left(
\frac{\pi^2}{L^2}-k
\right).
\]
但 $L>\pi/\sqrt k$ 使括号为负，矛盾。因此任何两点距离都不超过 $\pi/\sqrt k$，得到\eqnrefs{eq:math-dg-bonnet-myers-diameter}。

半径 $R$ 的球面 $\mathbb S^n(R)$：截面曲率 $K=1/R^2$，Ricci 曲率 $\operatorname{Ric}=(n-1)/R^2$，Bonnet--Myers 给出直径 $\le\pi R$，恰好等于球面直径——球面是等号情形。平坦环面 $T^n$：$\operatorname{Ric}=0$，定理不适用（$k=0$），直径确实无上界。双曲空间 $\mathbb H^n$：$\operatorname{Ric}=-(n-1)<0$，定理也不适用，测地线无限延伸。正曲率流形有限直径是整体微分几何的核心结论：它把局部曲率信息直接转化为全局紧致性，是 Cheeger--Gromoll 理论的起点。
\end{derivation}

这一证明揭示了 比较 几何 的主线：曲率下界本身并不直接``看见''全局直径；它先进入二阶变分的 指标 形式，再通过所有横向变分方向求和变成 Ricci 曲率，最后与最短性矛盾。局部张量信息由此转化为全局拓扑与度量结论。


Bonnet--Myers 使用的是固定端点二阶变分；体积比较则更适合把 Jacobi 方程改写成一阶矩阵 Riccati 方程。仍固定基点 $p$，在 $M\setminus(\{p\}\cup\operatorname{Cut}(p))$ 上记
\[
r(x)=d(p,x),
\qquad
\partial_r:=\nabla r.
\]
这里把函数上的几何 Laplace--Beltrami 算子显式记为
\begin{equation*}\Delta_{\mathrm{LB}}f:=\operatorname{tr}_g(\nabla^2f)=\nabla^i\nabla_i f
=-\Delta_H f.\end{equation*}
沿一条径向单位速测地线 $\gamma$，在法向空间 $T^\perp\gamma$ 上定义形算子
\begin{equation*}A_r(X):=\nabla_X\partial_r.\end{equation*}
它就是测地球面 $\{r=\text{const}\}$ 的第二基本形式所对应的自伴算子。若 $J$ 是满足 $J(0)=0$ 的横向 Jacobi 场，则在没有共轭点的区间内可写成
\[
J'(r)=A_rJ(r).
\]
把它代入 Jacobi 方程 $J''+R(J,\partial_r)\partial_r=0$，得到矩阵 Riccati 方程
\begin{equation*}\boxed{
A_r'+A_r^2+R_{\partial_r}=0,
\qquad
R_{\partial_r}(X):=R(X,\partial_r)\partial_r.}\end{equation*}
这一步把“二阶 Jacobi 方程”压缩成“测地球面形状的一阶演化方程”。

取迹并记
\[
H(r):=\tr A_r=\Delta_{\mathrm{LB}} r,
\]
便得到
\begin{equation*}H'(r)+|A_r|^2+\operatorname{Ric}(\partial_r,\partial_r)=0.\end{equation*}
若 $\operatorname{Ric}\ge(n-1)k g$，由 Cauchy--Schwarz
\[
|A_r|^2\ge\frac{H(r)^2}{n-1},
\]
于是
\begin{equation*}H'(r)+\frac{H(r)^2}{n-1}+(n-1)k\le0.\end{equation*}
模型函数 $S_k$ 已在\eqnrefs{eq:math-dg-model-jacobi}定义。令
\[
m_k(r):=(n-1)\frac{S_k'(r)}{S_k(r)},
\]
则 $m_k$ 满足相应的等号 Riccati 方程，并且 $r\downarrow0$ 时
\[
H(r)\sim\frac{n-1}{r},
\qquad
m_k(r)\sim\frac{n-1}{r}.
\]
标量 Riccati 比较因此给出 拉普拉斯算子 比较：
\begin{equation}
\boxed{
\Delta_{\mathrm{LB}} r\le(n-1)\frac{S_k'(r)}{S_k(r)}
}
\label{eq:math-dg-laplacian-comparison}
\end{equation}
在 割迹 之外逐点成立；在更一般的全局表述中可理解为分布或 barrier 意义的不等式。$k=0$ 时它退化为
\[
\Delta_{\mathrm{LB}} r\le\frac{n-1}{r},
\]
也就是说 Ricci 非负会使测地球面的平均曲率不超过 Euclidean 模型。

\begin{derivation}[从 Jacobi 行列式到 Bishop--Gromov 单调性]
在指数映射的极坐标中，体积元可写成
\[
\operatorname{vol}_g=J(r,\theta)\,\dd r\,\dd\theta,
\]
其中 $J(r,\theta)$ 是 $\dd\exp_p$ 在横向方向的 Jacobi 行列式。由行列式微分公式，
\[
\partial_r\log J(r,\theta)
=\tr A_r
=\Delta_{\mathrm{LB}} r.
\]
结合\eqnrefs{eq:math-dg-laplacian-comparison}，
\[
\partial_r
\log\frac{J(r,\theta)}{S_k(r)^{n-1}}
\le0.
\]
又因 $r\downarrow0$ 时
\[
J(r,\theta)\sim r^{n-1},
\qquad
S_k(r)^{n-1}\sim r^{n-1},
\]
所以对每个未越过割时刻的方向，
\begin{equation}
\frac{J(r,\theta)}{S_k(r)^{n-1}}
\quad\text{随 $r$ 单调不增。}
\label{eq:math-dg-jacobian-comparison}
\end{equation}
把横向角变量积分，再对半径积分，就得到 Bishop--Gromov 体积比较。若 $V_k(r)$ 表示常截面曲率 $k$ 的单连通模型空间中半径 $r$ 的球体积，则
\begin{equation}
\boxed{
\frac{\operatorname{Vol}_g(B_p(r))}{V_k(r)}
\quad\text{随 $r$ 单调不增。}
}
\label{eq:math-dg-bishop-gromov}
\end{equation}
特别地，当 $\operatorname{Ric}\ge0$ 时，
\[
\frac{\operatorname{Vol}_g(B_p(r))}{\omega_n r^n}
\]
单调不增且不超过 $1$。因此正或非负 Ricci 曲率不仅限制测地线聚焦，也限制体积增长速度。
\end{derivation}

横向 Jacobian 的单调性（\eqnrefs{eq:math-dg-jacobian-comparison}）积分即得 Bishop--Gromov 体积比较（\eqnrefs{eq:math-dg-bishop-gromov}）。

从逻辑上看，Rauch、Bonnet--Myers、拉普拉斯算子 比较 与 Bishop--Gromov 并不是四个彼此无关的定理：它们都从同一个 Jacobi 方程出发。Rauch 比较 Jacobi 场的长度；Bonnet--Myers 把 Jacobi 信息送入 指标 形式；拉普拉斯算子 比较 取 Riccati 方程的迹；Bishop--Gromov 再把迹积分为指数映射的 Jacobian。比较几何的“局部曲率 $\Rightarrow$ 全局度量结论”由此形成一条连续的推导链。

\begin{exbox}[极坐标中的平坦 Euclidean 平面]
在 $\RR^2\setminus\{0\}$ 上取
\[
g=\dd r^2+r^2\dd\phi^2,
\qquad
\theta^1=\dd r,
\qquad
\theta^2=r\dd\phi.
\]
这是一组正交归一余标架，但不是坐标余标架，因为
\[
\dd\theta^2
=\dd r\wedge\dd\phi
=\frac1r\theta^1\wedge\theta^2.
\]
无挠结构方程与\eqnrefs{eq:math-dg-orthonormal-connection-skew}给出
\[
\omega^1{}_2=-\frac1r\theta^2,
\qquad
\omega^2{}_1=\frac1r\theta^2.
\]
再代入第二结构方程，
\[
\Omega^1{}_2
=\dd\omega^1{}_2
+\omega^1{}_a\wedge\omega^a{}_2=0.
\]
联络一形式非零而曲率为零，说明“基底随位置转动”与“空间本身弯曲”是不同概念。
\end{exbox}
